\documentclass[10pt,a4paper]{amsart}
\usepackage{amsmath,amssymb,amsthm,graphicx}
% Retained preamble for the Stacks Project chapter snapshot in this directory.
% The chapter's own preamble.tex is a copy of the shared Stacks Project preamble; the block
% below is sliced MECHANICALLY, line for line, from it. Nothing was added, removed or
% redefined; the wrapper class, the xr-hyper/hyperref wiring and the externaldocument list
% are omitted because the assembler supplies the document class.
%
% lines taken from preamble.tex: 3, 5, 9, 11, 12, 156, 157, 158, ... (51 lines total)
%
% package coverage restored from the chapter's own preamble (2026-10-13): xy, xy-2cell

\usepackage[all]{xy}\usepackage{verbatim}
\xyoption{2cell}
\UseAllTwocells

\newenvironment{reference}{\comment}{\endcomment}
\newenvironment{slogan}{\comment}{\endcomment}
\newenvironment{history}{\comment}{\endcomment}
\theoremstyle{plain}
\newtheorem{theorem}[subsection]{Theorem}
\newtheorem{proposition}[subsection]{Proposition}
\newtheorem{lemma}[subsection]{Lemma}
\theoremstyle{definition}
\newtheorem{definition}[subsection]{Definition}
\newtheorem{example}[subsection]{Example}
\newtheorem{exercise}[subsection]{Exercise}
\newtheorem{situation}[subsection]{Situation}
\theoremstyle{remark}
\newtheorem{remark}[subsection]{Remark}
\newtheorem{remarks}[subsection]{Remarks}
\numberwithin{equation}{subsection}
\def\lim{\mathop{\mathrm{lim}}\nolimits}
\def\colim{\mathop{\mathrm{colim}}\nolimits}
\def\Spec{\mathop{\mathrm{Spec}}}
\def\Hom{\mathop{\mathrm{Hom}}\nolimits}
\def\Ext{\mathop{\mathrm{Ext}}\nolimits}
\def\SheafHom{\mathop{\mathcal{H}\!\mathit{om}}\nolimits}
\def\SheafExt{\mathop{\mathcal{E}\!\mathit{xt}}\nolimits}
\def\Sch{\mathit{Sch}}
\def\Mor{\mathop{\mathrm{Mor}}\nolimits}
\def\Ob{\mathop{\mathrm{Ob}}\nolimits}
\def\Sh{\mathop{\mathit{Sh}}\nolimits}
\def\NL{\mathop{N\!L}\nolimits}
\def\CH{\mathop{\mathrm{CH}}\nolimits}
\def\proetale{{pro\text{-}\acute{e}tale}}
\def\etale{{\acute{e}tale}}
\def\QCoh{\mathit{QCoh}}
\def\Ker{\mathop{\mathrm{Ker}}}
\def\Im{\mathop{\mathrm{Im}}}
\def\Coker{\mathop{\mathrm{Coker}}}
\def\Coim{\mathop{\mathrm{Coim}}}
\DeclareMathSymbol{\boxtimes}{\mathbin}{AMSa}{"02}
\def\QCohstack{\mathcal{QC}\!\mathit{oh}}
\def\Cohstack{\mathcal{C}\!\mathit{oh}}
\def\Spacesstack{\mathcal{S}\!\mathit{paces}}
\def\Quotfunctor{\mathrm{Quot}}
\def\Hilbfunctor{\mathrm{Hilb}}
\def\Curvesstack{\mathcal{C}\!\mathit{urves}}
\def\Polarizedstack{\mathcal{P}\!\mathit{olarized}}
\def\Complexesstack{\mathcal{C}\!\mathit{omplexes}}
\def\Pic{\mathop{\mathrm{Pic}}\nolimits}
\def\Picardstack{\mathcal{P}\!\mathit{ic}}
\def\Picardfunctor{\mathrm{Pic}}
\def\Deformationcategory{\mathcal{D}\!\mathit{ef}}
\usepackage{cleveref}
\crefname{originallemma}{original lemma}{original lemmas}
\Crefname{originallemma}{Original lemma}{Original lemmas}
\makeatletter
\newlabel{lemma-etale-local-finite-inertia}{{lemma-etale-local-finite-inertia}{source}{}{sourcepointer-1}{}}
\newlabel{lemma-etale-local-finite-inertia@cref}{{[originallemma][0][]lemma-etale-local-finite-inertia}{[1][0][]source}{}{}{}}
\makeatother
\crefname{originallemma}{original lemma}{original lemmas}
\Crefname{originallemma}{Original lemma}{Original lemmas}
\makeatletter
\newlabel{lemma-well-nigh-affine-moduli-space}{{lemma-well-nigh-affine-moduli-space}{source}{}{sourcepointer-2}{}}
\newlabel{lemma-well-nigh-affine-moduli-space@cref}{{[originallemma][0][]lemma-well-nigh-affine-moduli-space}{[1][0][]source}{}{}{}}
\makeatother
\crefname{originallemma}{original lemma}{original lemmas}
\Crefname{originallemma}{Original lemma}{Original lemmas}
\makeatletter
\newlabel{stacks-morphisms-lemma-aut-iso-unramified}{{0DU9}{source}{}{sourcepointer-3}{}}
\newlabel{stacks-morphisms-lemma-aut-iso-unramified@cref}{{[originallemma][0][]0DU9}{[1][0][]source}{}{}{}}
\makeatother
\crefname{originallemma}{original lemma}{original lemmas}
\Crefname{originallemma}{Original lemma}{Original lemmas}
\makeatletter
\newlabel{stacks-morphisms-remark-identify-automorphism-groups}{{0DTW}{source}{}{sourcepointer-4}{}}
\newlabel{stacks-morphisms-remark-identify-automorphism-groups@cref}{{[originallemma][0][]0DTW}{[1][0][]source}{}{}{}}
\makeatother
\crefname{originallemma}{original lemma}{original lemmas}
\Crefname{originallemma}{Original lemma}{Original lemmas}
\makeatletter
\newlabel{theorem-keel-mori}{{theorem-keel-mori}{source}{}{sourcepointer-5}{}}
\newlabel{theorem-keel-mori@cref}{{[originallemma][0][]theorem-keel-mori}{[1][0][]source}{}{}{}}
\makeatother
\title{Etale separated morphisms commute with the Keel-Mori moduli space}
\author{The Stacks Project Authors}
\date{stacks-project:stacks-more-morphisms.tex:local-snapshot-20261012; GNU FDL 1.2 or later}
\begin{document}
\maketitle
\noindent\textit{Editorial note.} Statement span is lines 4107-4129 of sources/stacks-stacks-more-morphisms/pinned-original-source.tex; proof span is lines 4131-4189, the complete author proof (57 body lines).\par\medskip
\noindent\textit{Editorial note.} Navigation stubs: 5 cited-but-not-included labels are supplied as non-load-bearing pointers with their pinned tags where the snapshot tag index has them.\par\medskip
\noindent\textit{Editorial note.} extra\_packages is empty on purpose: the retained preamble already loads xy with options and a second declaration clashes in the assembler wrapper.\par\medskip
\noindent\textit{Editorial note.} Text is retained verbatim, including the author's own typos and placeholders.\par\medskip
\noindent\textit{Editorial note.} Queue index 23 of queue\_stacks-stacks-more-morphisms.json (pb=2178, stmt 4107-4129, proof 4131-4189).\par\medskip
\noindent\textit{Editorial note.} stacks-stacks-more-morphisms wave smm01 (2026-10-14): fresh adjudication accept per lane-s3/smm\_judgments.md; pool per Chief 10:5x open-pool ruling (net 59 with-proof rows, bloodline lane-s3 by anchor audit, v2 true-key namespace stacks:stacks-more-morphisms.tex, row lemma-deform-property-fp-over-ft avoided on cross-chapter key collision); 60 with-proof rows judged 45 accept / 14 reject / 1 avoid; patch-sd-envs contract spans scanned.\par\medskip
\noindent\textit{Allowed standard prerequisites (editorial).}\begin{itemize}
\item Explicit prerequisite boundary. The cited inputs are prerequisites; the argument in the retained span is not replaced by any cited result.
\end{itemize}\medskip
\ifdefined\hypertarget\hypertarget{sourcepointer-1}{}\fi
\noindent\textit{Source reference.} Original lemma lemma-etale-local-finite-inertia: sources/stacks-stacks-more-morphisms/pinned-original-source.tex, source lines 3798--3814. Cited by the retained statement or proof but not part of the collected spans; supplied only as a navigation pointer. The argument does not depend on its content.\par\medskip
\ifdefined\hypertarget\hypertarget{sourcepointer-2}{}\fi
\noindent\textit{Source reference.} Original lemma lemma-well-nigh-affine-moduli-space: sources/stacks-stacks-more-morphisms/pinned-original-source.tex, source lines 3425--3433. Cited by the retained statement or proof but not part of the collected spans; supplied only as a navigation pointer. The argument does not depend on its content.\par\medskip
\ifdefined\hypertarget\hypertarget{sourcepointer-3}{}\fi
\noindent\textit{Source reference.} Original lemma 0DU9: sources/stacks-stacks-more-morphisms/pinned-original-source.tex, source lines 4189--4189. Cited by the retained statement or proof but not part of the collected spans; supplied only as a navigation pointer. The argument does not depend on its content.\par\medskip
\ifdefined\hypertarget\hypertarget{sourcepointer-4}{}\fi
\noindent\textit{Source reference.} Original lemma 0DTW: sources/stacks-stacks-more-morphisms/pinned-original-source.tex, source lines 4189--4189. Cited by the retained statement or proof but not part of the collected spans; supplied only as a navigation pointer. The argument does not depend on its content.\par\medskip
\ifdefined\hypertarget\hypertarget{sourcepointer-5}{}\fi
\noindent\textit{Source reference.} Original lemma theorem-keel-mori: sources/stacks-stacks-more-morphisms/pinned-original-source.tex, source lines 3854--3862. Cited by the retained statement or proof but not part of the collected spans; supplied only as a navigation pointer. The argument does not depend on its content.\par\medskip
\begin{lemma}
\label{lemma-etale-separated-over-keel-mori}
Let $h : \mathcal{X}' \to \mathcal{X}$ be a morphism of algebraic stacks.
Assume
\begin{enumerate}
\item $\mathcal{I}_\mathcal{X} \to \mathcal{X}$ is finite,
\item $h$ is \'etale, separated, and induces isomorphisms on
automorphism groups (Morphisms of Stacks, Remark
\ref{stacks-morphisms-remark-identify-automorphism-groups}).
\end{enumerate}
Then there exists a cartesian diagram
$$
\xymatrix{
\mathcal{X}' \ar[d] \ar[r] &
\mathcal{X} \ar[d] \\
M' \ar[r] &
M
}
$$
where $M' \to M$ is a separated \'etale morphism of algebraic spaces and
the vertical arrows are the moduli spaces constructed in
Theorem \ref{theorem-keel-mori}.
\end{lemma}

\begin{proof}
By Morphisms of Stacks, Lemma \ref{stacks-morphisms-lemma-aut-iso-unramified}
we see that
$\mathcal{I}_{\mathcal{X}'} \to
\mathcal{X}' \times_\mathcal{X} \mathcal{I}_\mathcal{X}$
is an isomorphism. Hence $\mathcal{I}_{\mathcal{X}'} \to \mathcal{X}'$
is finite as a base change of $\mathcal{I}_\mathcal{X} \to \mathcal{X}$.
Let $f' : \mathcal{X}' \to M'$ and $f : \mathcal{X} \to M$ be as in
Theorem \ref{theorem-keel-mori}.
We obtain a commutative diagram as in the lemma because
$f'$ is categorical moduli space.
Choose $I$ and $g'_i : \mathcal{X}'_i \to \mathcal{X}'$ as in
Lemma \ref{lemma-etale-local-finite-inertia}.
Observe that $g_i = h \circ g'_i$
is \'etale, separated, and induces isomorphisms on
automorphism groups (Morphisms of Stacks, Remark
\ref{stacks-morphisms-remark-identify-automorphism-groups}).
Let $f'_i : \mathcal{X}'_i \to M'_i$ be as in
Lemma \ref{lemma-well-nigh-affine-moduli-space}.
In the proof of Theorem \ref{theorem-keel-mori}
we have seen that the diagrams
$$
\xymatrix{
\mathcal{X}'_i \ar[d]_{f'_i} \ar[r]_{g'_i} &
\mathcal{X}' \ar[d]^{f'} \\
M'_i \ar[r] &
M'
}
\quad\text{and}\quad
\xymatrix{
\mathcal{X}'_i \ar[d]_{f'_i} \ar[r]_{g_i} &
\mathcal{X} \ar[d]^f \\
M'_i \ar[r] &
M
}
$$
are cartesian and that $M'_i \to M'$ and $M'_i \to M$ are \'etale
(this also follows directly from the properties of the morphisms
$g'_i, g_i, f', f'_i, f$ listed sofar by arguing in exactly the same way).
This first implies that $M' \to M$ is \'etale and second that
the diagram in the lemma is cartesian. We still need to show
that $M' \to M$ is separated. To do this we contemplate the
diagram
$$
\xymatrix{
\mathcal{X}' \ar[r] \ar[d] &
\mathcal{X}' \times_\mathcal{X} \mathcal{X}' \ar[d] \\
M' \ar[r] &
M' \times_M M'
}
$$
The top horizontal arrow is universally closed as
$\mathcal{X}' \to \mathcal{X}$ is separated.
The vertical arrows are as in Theorem \ref{theorem-keel-mori}
(as flat base changes of $\mathcal{X} \to M$)
hence universal homeomorphisms. Thus the lower horizontal
arrow is universally closed. This (combined with it being an \'etale
monomorphism of algebraic spaces) proves it is a closed immersion as desired.
\end{proof}

\end{document}
